\chapter{المكابح}
% full version: chapters/03_gaba.tex

لنضف إلى الآلة ثاني أكثر أنواع العصبونات عددًا: \nt{GABA}. إشارتها لا تدفع الجار نحو النقر، بل تبعده عنه. إنها المكابح. عددها قليل، في القشرة من الخُمس إلى الربع من كل العصبونات، لكن الآلة لا تسير من دونها.

ثمة قيد واحد: العصبون ”الزائد“ لا يستطيع أن يكبح جاره بنفسه. إذا احتجنا إلى وصلة سالبة، وجب بناؤها عبر وسيط: ”الزائد“ يوقظ عصبونًا كابحًا، وهذا بدوره يكبح الهدف. قلب ”الزائد“ إلى ”ناقص“ يكلف عصبونًا واحدًا وتأخيرًا صغيرًا واحدًا.

\section*{”إلا إذا“}

أنفع عملية بالمكابح هي المنع: ”اعمل عند وصول أ، إلا إذا وصل ب“. يتلقى العصبون دفعة من أ، وكبحًا من وسيط يوقظه ب. ومن المنع و”أو“ يمكن تجميع أي منطق.

\pencilfig{3}{\pencilart{3}}{”أ، إلا إذا ب“: العصبون الأحمر يسد الطريق على الأزرق.}

المنع في الزمن يعطي كاشف الحركة الموجود في عينك. مستشعران متجاوران: أحدهما يثير عصبون الخرج، والآخر يكبحه بتأخير. إذا جرت بقعة ضوء في اتجاه ما بسرعة مناسبة، وصل الكبح في لحظة وصول الإثارة نفسها، فأطفأها. وإذا جرت في الاتجاه الآخر، تأخر الكبح، ولحق العصبون أن ينقر. فنحصل على عصبون لا يرى الحركة إلا في اتجاه واحد.

\pencilfig{103}{%
% sensors on top; the left one excites the output, the right one inhibits it through a GABA go-between and a delay
\draw[dotpen=pcAmber, wash=pcAmber] (0,1.6) circle (0.16); \draw[dotpen=pcAmber, wash=pcAmber] (1.6,1.6) circle (0.16);
\node[lbl, anchor=south] at (0.8,1.8) {مستشعران};
\draw[pen=pcBlue, wash=pcBlue] (0.4,0) circle (0.24);
\draw[arr=pcBlue] (0,1.44) -- (0.3,0.25);
\draw[pen=pcBlue] (1.6,1.44) -- (1.6,1.18);
\draw[penlite, fill=white] (0.95,0.9) rectangle (2.25,1.18); \node[font=\tiny\itshape, text=pcGray] at (1.6,1.04) {تأخير};
\draw[pen=pcBlue] (1.6,0.9) -- (1.6,0.72);
\draw[pen=pcRed, wash=pcRed] (1.6,0.55) circle (0.17);
\draw[pcRed, line width=0.7pt, -{Bar[width=5pt]}] (1.45,0.45) -- (0.66,0.08);
\draw[arr] (0.4,-0.24) -- (0.4,-0.6); \node[lbl, anchor=north] at (0.4,-0.6) {الخرج};
% outcomes
\begin{scope}[xshift=3.1cm]
  \draw[dotpen=pcAmber, wash=pcAmber] (0,1.25) circle (0.1); \draw[arr=pcAmber] (0.18,1.25) -- (1.1,1.25);
  \node[lbl, anchor=west] at (1.25,1.25) {من اليسار إلى اليمين: نقرة};
  \draw[dotpen=pcAmber, wash=pcAmber] (1.1,0.45) circle (0.1); \draw[arr=pcAmber] (0.92,0.45) -- (0,0.45);
  \node[lbl, anchor=west] at (1.25,0.45) {من اليمين إلى اليسار: صمت};
  \node[lbl, anchor=west, align=left] at (0,-0.35) {يصل الكبح\\مع الإثارة تمامًا};
\end{scope}
}{كاشف الحركة: يصل الكبح متأخرًا، فلا يطفئ الاستجابة إلا إذا جرت البقعة من اليمين إلى اليسار.}


\section*{طرح أم قسمة}

لا يقتصر الكبح على الطرح. القناة الكابحة المفتوحة لا تشد الجهد إلى الأسفل، بل تشده \emph{نحو الراحة}، كأننا ثقبنا الدلو ثقبًا آخر. فيصير أي تدفق للماء يرفعه أقل. الكبح لا ينتزع من الإشارة حصة ثابتة، بل \emph{يقسمها}: إنه مقبض لضبط الصوت. وإذا نمت قوة الكبح مع النشاط العام من حوله، استجاب كل عصبون لا لقوة إشارته، بل لقوتها \emph{مقارنة بجيرانه}. لذلك تعمل الشبكة نفسها في الغسق وتحت الشمس الساطعة. لكن هذا ”المقسِّم“ لا يؤثر في تردد النبضات تأثيرًا نظيفًا إلا مع ضجيج خلفي ثابت، وهو موجود دائمًا في الدماغ الحي.

\pencilfig{104}{%
\foreach \dx/\t in {0/طرح,4.6/قسمة}{
  \begin{scope}[xshift=\dx cm]
  \draw[pen] (0,0) -- (3.6,0); \draw[pen] (0,0) -- (0,1.9);
  \node[lbl, anchor=north] at (1.8,-0.05) {الإشارة}; \node[lbl, anchor=south, rotate=90] at (-0.05,0.95) {الاستجابة};
  \draw[pen=pcBlue] (0.4,0) -- (3.3,1.75);
  \node[font=\small\itshape, text=pcGray, anchor=south] at (1.8,1.95) {\t};
  \end{scope}}
\draw[pen=pcRed, line width=0.9pt] (1.3,0) -- (3.4,1.27);
\node[lbl, text=pcRed, anchor=west] at (2.25,0.35) {إزاحة};
\begin{scope}[xshift=4.6cm]
\draw[pen=pcRed, line width=0.9pt] (0.4,0) -- (3.3,0.8);
\node[lbl, text=pcRed, anchor=west] at (2.0,0.25) {ميل أقل};
\end{scope}
\node[lbl, text=pcBlue, anchor=east] at (2.25,1.45) {بلا كبح};
}{بالأزرق استجابة العصبون بلا كبح، وبالأحمر معه. الطرح يزيح العتبة؛ والقسمة تخفض الصوت، كمقبض الضبط.}

وثمة نتيجة ثانية: الثقب الإضافي يجعل الدلو ينسى أسرع. فتضيق النافذة التي يجب أن تقع فيها إشارتان كي تُجمعا. الكبح يجعل كواشف التزامن أكثر صرامة.

\section*{الاختيار}

والآن أهم ما كان ينقصنا: اختيار واحد من كثير. مجموعتان من العصبونات تكبح كل منهما الأخرى. ما دام الكبح ضعيفًا، تعمل الاثنتان، لكن بصوت أخفض. أما إذا كان قويًا بما يكفي، فأي فرق صغير يتضخم، كما في الأرجوحة، حتى تصمت إحدى المجموعتين تمامًا. الفائز يأخذ كل شيء. ويثبت الفوز: كي تغير الشبكة رأيها، يجب أن تكون الحجة الجديدة أقوى من القديمة بفارق ملحوظ. فلا ترتجف الآلة مع كل حفيف.

والمكابح تمحو الذاكرة أيضًا: عصبون كابح يطفئ، عند إشارة ”إعادة الضبط“، نار المخيم من الفصل السابق. وخليتان كهاتين تكبح كل منهما الأخرى هما نظير المِزلاج، أساس كل ذاكرة حاسوبية.

\section*{حياة على الحافة}

شبكة من ”الزوائد“ وحدها قد تنزلق إلى انهيار متتابع: كل نبضة تولد أكثر من نبضة تالية. المكابح تبقيها عند نقطة العمل، كما يبقي منظم الحرارة حرارة الفرن. ولذلك شرطان: أن يكون الكبح قويًا بما يكفي وسريعًا بما يكفي. فإذا كان قويًا لكنه بطيء، أخذ المنظم يبالغ في التصحيح، كمن يدير صنبور الدش دون أن ينتظر وصول الماء الساخن. فتتأرجح الحرارة ذهابًا وإيابًا. وفي الدماغ تولّد حلقة كهذه من ”الزوائد“ و”النواقص“ إيقاعًا سريعًا: بضع عشرات من الذبذبات في الثانية.

ويبدو أن القشرة تعيش على هذه الحافة بالذات: تغذيتها الراجعة الذاتية قوية بما يكفي لتنزلق إلى الانهيار لولا المكابح. لا يمسكها إلا الكبح السريع، ولهذا تستجيب بهذه الحساسية للإشارات الضعيفة.

خلاصة الفصل غير متوقعة: المكابح لا تضيف إلى الآلة حسابات جديدة تقريبًا. إنها تضيف \emph{التحكم}: الاختيار، وإعادة الضبط، وضبط الصوت، والاستقرار. الإثارة تحمل البيانات، والكبح يقرر أي البيانات تمر، ومتى، وبأي علوّ.

\fullversion{3}
